\section{K-036: a uniform criterion excluding rank-two odd inertia}

The positive-dimensional step of K-035 has a geometric version.
It does not require an enumeration of balance vertices and applies
to arbitrary numbers of frame vectors.  Its hypotheses concern
the full stabilizer, the complete fixed support, and the actual
critical masses of a smaller tensor; none is inferred from an
incidence determinant alone.

\subsection*{The hypotheses and the mass condition}

Let $m\ge3$, and let $T_0\in(\mathbb C^m)^{\otimes3}$ be a critical tensor on a free
support $\Sigma$ with $3m-2$ independent incidence columns and
positive margin-one masses $p_e$.  Suppose its full stabilizer is
$S_m\times\mu_D$, with faithful determinant-product character,
where $D$ is a power of two.  Suppose a generator $h_0$ of its
finite factor has simple spectra, complete zero-weight support
$\Sigma$, and antipodal-free spectra in two of the factors.

Fix $e_0\in\Sigma$ of mass $\rho$, with
\[
                              \frac12<\rho<1,
                  \qquad t=1-\rho.                     \tag{K-036a}
\]
Let $\mathcal E_a$ be the old edges other than $e_0$ incident to
its selected coordinate in factor $a$.  Freeness makes these
three edge sets disjoint.  Each is nonempty and its masses sum
to $t$.  Define its finite set of subset sums by
\[
   F_a=\left\{\sum_{e\in I}p_e:I\subseteq\mathcal E_a\right\}.
\]
Assume the following condition for each pair of distinct factors:
\[
                             F_a\cap F_b=\{0,t\}.       \tag{K-036b}
\]
It includes complementary allocations, since $F_b=t-F_b$.
When one factor has a single frame edge and the other two have
two each, only the two latter sets need checking.  The condition
then says that their two unordered mass pairs have no common entry.

Duplicate the selected coordinate in each factor and extend
$h_0$ by the same eigenvalue on the new coordinate.  Denote this
order-$D$ element by $h$, its complete fixed space by $U$, and
the reduced closed image of $U$ in $X_{m+1}$ by $Y$.
Every old edge group other than $e_0$ is either a scalar or a
two-vector; $e_0$ becomes the full binary cube.

\begin{theorem}[K-036]
Every nonzero point of $Y$ has grading stabilizer $\mu_D$ or
$\mu_{2D}$.  In particular $Y$ has no nonzero point with a grading
stabilizer of order divisible by an odd prime.

If, in addition, $e_0$ has the pair-deletion coverage of K-015
and $\rho$ has reduced denominator $D$, then $Y$ is an actual
irreducible component of $X_{m+1}^{\mu_D}$.  It contains an
effectively free smooth quotient point and has dimension
\[
                       |\mathcal E_1|+|\mathcal E_2|
                                  +|\mathcal E_3|-1.    \tag{K-036c}
\]
\end{theorem}

\subsection*{Critical tensors with a positive-dimensional centralizer}

We first prove the only new geometric step.  Choose a critical
representative $T\in U$ and normalize its reduced matrices to
identity.  The existence of such a representative over every
nonzero point of $Y$ follows from the finite restriction and
moment-map argument of K-032; that argument uses equality of
nullcones, not generic freeness.  In particular its image is
closed.  The aggregate group masses are $p_e$ for $e\ne e_0$
and $1+\rho=2-t$ for the core.  Every frame is nonzero.

As in K-035, the full stabilizer of a critical tensor is
adjoint-closed.  Each selected two-dimensional factor has a
nonzero invariant frame line under the centralizer of $h$ in
that stabilizer.  Its orthogonal complement is invariant too.
A block-unitary basis change therefore diagonalizes that entire
centralizer.  If its quotient by the scalar torus is
positive-dimensional, choose a non-scalar cocharacter in it.
Let $\delta_a$ be the difference between its two weights in
the selected block of factor $a$.

The vector $\delta=(\delta_1,\delta_2,\delta_3)$ is not zero.
Otherwise the weights collapse to old-coordinate weights fixing
every old group; independence of the old incidence columns
would make the cocharacter scalar.  On the core, supported cube
vertices $b\in\{0,1\}^3$ satisfy
\[
                              \delta\cdot b=c           \tag{K-036d}
\]
for one constant $c$.  Thus their affine span has dimension at
most two.

Let $f_a$ be the total frame mass in the coordinate labelled
one in factor $a$.  Always $0\le f_a\le t$.  If
$\delta_a\ne0$, each frame vector in that factor is an
eigenvector of the cocharacter and lies on one of the two
coordinate lines.  In that case, and only where needed below,
we know that $f_a\in F_a$.  The normalized barycenter of the
core support has coordinates
\[
              b_a^*=\frac{1-f_a}{2-t},\qquad
       \frac{1-t}{2-t}\le b_a^*\le\frac1{2-t}.          \tag{K-036e}
\]
The same interval is preserved by reflecting any cube coordinate.
All its entries are strictly between zero and one.  Moreover,
because $t<1/2$, every point in this box has coordinate sum
strictly between one and two.

Suppose first that the affine span of the supported cube
vertices has dimension two.  Choose three noncollinear such
vertices and reflect coordinates so that the first is the
origin.  The other two are binary vectors.  Their nonzero cross
product has coordinates in $\{-1,0,1\}$ and is a normal vector
to the plane.  A normal with exactly one nonzero coordinate
would make one barycenter coordinate zero or one, which is
impossible.  If all three coordinates were nonzero, further
reflections would put the plane in the form
$b_1+b_2+b_3=1$ or $2$: these are the only integer-level
planes of that form containing three noncollinear cube vertices.
This contradicts the last observation after (K-036e).

The normal therefore has exactly two nonzero coordinates.  After
reflections the plane is $b_a=b_b$ for some $a\ne b$.  These
are precisely the two nonzero coordinates of $\delta$.
The corresponding reflected frame totals satisfy $f_a=f_b$,
and both are subset sums in their respective $F$ sets.
Condition (K-036b) forces them to be zero or $t$.  Thus all
frames in those two factors lie on the same side of the
two-coordinate diagonal.  The other coordinate lines are empty
of frames and are the two isolated lines.  The whole tensor
is invariant under weight one on the first isolated line,
weight minus one on the second, and weight zero elsewhere.
The reduced-matrix argument (K-035i) now gives an orthogonal
one-dimensional tensor summand.  Indeed the first two isolated
margins force its remaining vector to have norm one; positivity
of the remaining reduced matrix forces the complementary tensor
to be orthogonal to that vector.

Suppose next that the supported cube vertices have affine span
one.  A line contains at most two distinct cube vertices.  The
strictly interior coordinates in (K-036e) imply that these two
vertices differ in every coordinate, so they are antipodal.
At least two coordinates of $\delta$ are nonzero: a single
nonzero coordinate would make that core coordinate constant.
Reflect the cube so that its segment joins $000$ and $111$.
All three barycenter coordinates are then equal.  For two
factors with nonzero $\delta_a$, the corresponding frame
totals are actual subset sums and are equal.  Condition
(K-036b) again makes their common value zero or $t$.
The two core masses are consequently $1$ and $1-t=\rho$.
In every factor the sum of the positive frame outer products
is the reduced matrix $\operatorname{diag}(0,t)$, up to
exchanging the two coordinates.  Every frame, including those
in a factor where $\delta_a=0$, therefore lies on the
$\rho$ side.  The mass-one core term is an orthogonal
one-dimensional tensor summand.

Affine span zero is impossible by the strict interior
barycenter condition.  We have covered all cases.  After
removing the one-dimensional summand, the remaining tensor
has the old support and exactly its positive mass vector.
Independence of the old incidence columns makes its nonzero
coefficient torus a single diagonal orbit.  Thus
\[
  \dim(C_{\operatorname{Stab}(T)}(h)/S)>0
           \quad\Longrightarrow\quad T\simeq T_0\oplus T_1.
                                                               \tag{K-036f}
\]
This proof also explains why complementary frame allocations
must be excluded in (K-036b), and why it is insufficient to
treat only a full-dimensional core support.

\subsection*{From the centralizer to the full stabilizer}

All remaining steps of K-035 now apply uniformly.  We spell out
their hypotheses.  The full stabilizer is reductive, contains
the scalar torus, and has finite determinant character.  The
first and third factors, after renumbering, have antipodal-free
$h$ spectra.  The stable maximal torus and faithful two-factor
representation argument of Lemma K-035b therefore shows that
a positive-dimensional full effective stabilizer has a
positive-dimensional effective $h$ centralizer.  By (K-036f)
the tensor is then $T_0\oplus T_1$.

The old full stabilizer hypothesis implies its Lie stabilizer
is $S_m$.  The centroid argument in K-035 identifies the
centroid of $T_0\oplus T_1$ with $\mathbb C\oplus\mathbb C$.
The ranks $m$ and one differ, so its full stabilizer preserves
the two summands.  Its determinant character image is exactly
$\mu_D$.

If the full effective stabilizer is finite, its $h$ centralizer
is diagonal after the frame-line argument.  Choose one
supported coordinate in each old group and extend the
$3m-2$ independent base columns by three columns.  The same
unimodular row transformation as in (K-035e) gives determinant
$\pm D\det B$, where $B$ is, up to row and column signs,
a binary $3\times3$ matrix.  The two primitive old block-sum row
relations have complementary $2\times2$ minors zero or $\pm1$.
Thus all nonzero old maximal incidence minors have the same
absolute value, equal to their greatest common divisor, the
support-lattice index.  That index is $D$, because the old
complete diagonal stabilizer is exactly $S_m\times\mu_D$.
Hence the finite
effective centralizer has order at most $2D$.

The effective centralizer equals the actual centralizer modulo
$S$: a scalar spectral shift must preserve the unique double
eigenvalue in every factor and is therefore trivial.  The
elementary conjugacy-coset inequality (K-035g) bounds the
order of the full determinant character image by $2D$.
That image contains $\mu_D$, because duplicating the selected
three eigenvalues multiplies their determinant product by
their product, which is one.  It is thus $\mu_D$ or
$\mu_{2D}$.  Identification with the quotient grading
stabilizer proves the first assertion of the theorem.

For the additional component assertion, perform the
edge-doubling construction K-015.  Pair-deletion coverage
gives Lie stabilizer $S$, and the odd numerator of $\rho$
at denominator $D$ gives the cyclic diagonal factor of
order $2D$ with faithful determinant character.  The
old simple spectra and complete support imply that a
primitive new generator has simple spectra with exactly
one antipodal pair in each factor and exactly the
constructed support.  Each nonempty frame set supplies
a nonzero one-paired-plane projection.  K-016 gives the
full stabilizer $S\times\mu_{2D}$.  The quotient point is
therefore effectively free and smooth.

Put $c=|\mathcal E_1|+|\mathcal E_2|+|\mathcal E_3|$.
The complete fixed tensor space has dimension
$(3m-2)+7+c$, and the fixed special-linear Lie algebra
has dimension $3(m+2)$.  The fixed tangent dimension is
$c-1$.  As in K-029, surjectivity onto this smooth fixed
germ identifies the irreducible fixed-space image with
the entire irreducible component.  This proves (K-036c).

\subsection*{Application to every amplified prefix}

Take the K-020 tensor after $k\ge1$ amplifications.  Then
\[
             m=2+7k,\qquad D=2\cdot256^k.
\]
Let $r\in I=(3/11,37/61)$ be the selected input mass
before its last amplification.  Its selected output
mass and the four two-frame masses are
\[
\begin{aligned}
 \rho&=(125+43r)/256,\\
 \alpha_1&=(-1+57r)/256,&
 \alpha_2&=(132-100r)/256,\\
 \gamma_1&=(136-72r)/256,&
 \gamma_2&=(-5+29r)/256.
\end{aligned}                                             \tag{K-036g}
\]
The middle factor has its single frame mass $t=1-\rho$.
All masses are positive by the K-018 interval certificate.
We have $\rho>1/2$ and
\[
          \gamma_2<\alpha_1<\alpha_2<\gamma_1
                         \qquad(r\in I),                \tag{K-036h}
\]
as follows directly from
$\alpha_1-\gamma_2=(4+28r)/256$,
$\alpha_2-\alpha_1=(133-157r)/256$, and
$\gamma_1-\alpha_2=(4+28r)/256$.
Thus (K-036b) holds.  All spectral, full-stabilizer,
coverage, and exact-denominator hypotheses were proved
uniformly in K-020.  The theorem gives, in every
dimension $n=3+7k$, an actual four-dimensional component
$Y_k\subset X_n^{\mu_D}$ with no nonzero odd grading inertia.

For the originally targeted order $M_n=2^{n-1}$,
\[
                            D=2^{k-1}M_n.               \tag{K-036i}
\]
Consequently $Y_k\subset X_n^{\mu_{M_n}}$ and
$Y_k\cap X_n^{\mu_{3M_n}}=\{0\}$.  When $k=1$ this
is precisely the whole component of K-035.  When
$k>1$, $Y_k$ may be a proper subvariety of a larger
$\mu_{M_n}$-fixed component.  The conclusion does not
remove that larger component from K-031's collision
support.  It proves an all-point inertia bound on the
specified entire $\mu_D$ component, not a global
noncancellation theorem at $M_n$.

\paragraph{Verification and scope.}
The new positive-centralizer argument above is a complete
cube-geometric proof.  It does not rely on the exploratory
hyperplane enumeration that suggested it.  The finite-group,
critical-point, stable-torus, and centroid inputs are
reconstructed in K-035, with its stated primary-source
scope.  Formula (K-036g) is the literal selection of
edges $755$, $746$, $770$, $275$, and $065$ from the
already checked K-018 gadget; its inequalities are
displayed explicitly.  The auxiliary verifier
\texttt{verify\_uniform\_rank\_two\_inertia.py} checks all
$56$ triples of cube vertices, all $28$ cube segments, and
the endpoint signs of every affine inequality in
(K-036g)--(K-036h); its exact receipt is
\texttt{uniform\_rank\_two\_inertia\_checks.json}.
No coefficient of a global
Hilbert series or strict tensor index is asserted here.
